% =====================================================================
%  Seasonal-Aware Economic Dispatch (SAED) and a Digital Twin
%  for the Bangladesh (PGCB) Power System
%  Compile with pdfLaTeX (Overleaf-ready). Requires the /figures and
%  /tables folders shipped alongside this file.
% =====================================================================
\documentclass[11pt]{article}

\usepackage[a4paper,margin=1in]{geometry}
\usepackage{amsmath,amssymb,amsthm}
\usepackage{booktabs}
\usepackage{graphicx}
\usepackage{subcaption}
\usepackage{caption}
\usepackage[ruled,vlined,linesnumbered]{algorithm2e}
\usepackage{siunitx}
\usepackage{xcolor}
\usepackage[hidelinks]{hyperref}
\usepackage{microtype}

\newtheorem{theorem}{Theorem}
\newtheorem{definition}{Definition}

\DeclareMathOperator*{\argmin}{arg\,min}
\newcommand{\Pmax}{P^{\mathrm{max}}}
\newcommand{\seasonset}{\mathcal{S}}
\newcommand{\srcset}{\mathcal{G}}

\title{\textbf{A Seasonal-Aware Economic Dispatch Algorithm and Digital Twin
for the Bangladesh Power System}\\[2pt]
\large Data-driven calibration, optimal merit-order dispatch,
and seasonal generation analysis}
\author{Generated analysis based on the PGCB hourly record (2015--2025)}
\date{\today}

\begin{document}
\maketitle

\begin{abstract}
\noindent
We develop a \emph{Seasonal-Aware Economic Dispatch} (SAED) algorithm and an
accompanying \emph{digital twin} for the Power Grid Company of Bangladesh
(PGCB) system, calibrated entirely from a $\sim$92{,}650-record hourly archive
spanning April 2015 to June 2025. The method (i) classifies each operating hour
into one of four Bangladeshi seasons, (ii) reconstructs a seasonal
demand model from a base level, a normalized diurnal profile and an annual
growth trend, (iii) estimates a per-source \emph{seasonal availability matrix}
that encodes fuel-supply, hydrological and maintenance limits as they actually
manifest, and (iv) solves a constrained economic-dispatch problem whose optimal
solution is the classical merit order. The digital twin replays the calibrated
model to produce publication-ready seasonal dispatch stacks, a validation
against the historical mix, and a demand-growth stress test. The twin reproduces
the qualitative seasonal physics of the grid---hydro peaking in the monsoon,
solar in the pre-monsoon, and liquid-fuel peaking concentrated in the
high-demand monsoon months---and quantifies the cost and carbon penalty of
seasonal peaking, together with the year in which each season's reserve margin
falls below the planning threshold.
\end{abstract}

% =====================================================================
\section{Dataset}
The model is calibrated on the PGCB hourly/half-hourly record with fields for
total \texttt{generation} and \texttt{demand} (MW), \texttt{load\_shedding}, and
the MW contribution of each supply source: domestic \texttt{gas},
\texttt{liquid\_fuel}, \texttt{coal}, \texttt{hydro}, \texttt{solar},
\texttt{wind}, and cross-border imports
(\texttt{India\,Bheramara\,HVDC}, \texttt{India\,Tripura}, \texttt{India\,Adani},
\texttt{Nepal}). Records whose values exceed a physical grid ceiling of
\SI{20}{GW} (data-entry artefacts) are discarded; the four import channels are
aggregated into a single dispatchable \texttt{imports} block.

% =====================================================================
\section{Mathematical formulation}

\subsection{Sets and seasonal classification}
Let the supply sources and seasons be
\begin{align*}
\srcset&=\{\text{solar},\text{hydro},\text{gas},\text{coal},
\text{imports},\text{liquid\_fuel}\},\\
\seasonset&=\{\text{Winter},\text{PreMonsoon},\text{Monsoon},
\text{PostMonsoon}\}.
\end{align*}
With $m\in\{1,\dots,12\}$ the calendar month, the season map
$\sigma:\{1,\dots,12\}\to\seasonset$ is
\begin{equation}
\sigma(m)=
\begin{cases}
\text{Winter}, & m\in\{12,1,2\},\\
\text{PreMonsoon}, & m\in\{3,4,5\},\\
\text{Monsoon}, & m\in\{6,7,8,9\},\\
\text{PostMonsoon}, & m\in\{10,11\}.
\end{cases}
\end{equation}

\subsection{Seasonal demand model}
Demand is decomposed into a seasonal base level $\bar{D}_s$, a normalized
diurnal shape $\phi_s(h)$, a multiplicative annual growth factor, and a
stochastic term:
\begin{equation}
D(s,h,y)=\bar{D}_s\,\phi_s(h)\,(1+g)^{\,y-y_0}\,\bigl(1+\varepsilon\bigr),
\qquad \varepsilon\sim\mathcal{N}\!\left(0,\eta_s^2\right),
\label{eq:demand}
\end{equation}
where $h\in\{0,\dots,23\}$ is the hour, $y$ the year, $y_0$ the base
(reference) year, $g$ the annual growth rate, and $\eta_s$ the season's relative
volatility. The diurnal profile is normalized so that it has unit mean,
\begin{equation}
\frac{1}{24}\sum_{h=0}^{23}\phi_s(h)=1
\quad\Longrightarrow\quad
\mathbb{E}_h\!\left[D(s,h,y)\right]=\bar{D}_s(1+g)^{y-y_0}.
\end{equation}

\subsection{Firm available capacity}
The firm power a source can deliver in season $s$ at hour $h$ is the product of
a firm rated capacity $C_g$, a dimensionless seasonal availability factor
$\alpha_{g,s}\in[0,1]$, and an intraday availability shape $\psi_g(h)$:
\begin{equation}
\Pmax_{g}(s,h)=C_g\,\alpha_{g,s}\,\psi_g(h).
\label{eq:pmax}
\end{equation}
For all dispatchable thermal/hydro/import sources $\psi_g(h)\equiv1$; only
solar carries a non-trivial daylight shape, $\psi_{\text{solar}}(h)\in[0,1]$,
with $\psi_{\text{solar}}=0$ overnight. The factor $\alpha_{g,s}$ is the
mechanism through which \emph{season changes the generation mix}: it captures
reservoir filling in the monsoon (high $\alpha_{\text{hydro}}$), clear-sky
irradiance in the pre-monsoon (high $\alpha_{\text{solar}}$), and the winter
contraction of domestic gas supply (low $\alpha_{\text{gas}}$).

\subsection{Economic dispatch as a linear program}
Let $c_g$ be the marginal cost (\si{USD/MWh}) of source $g$. For a single hour
with demand $D$, SAED solves
\begin{equation}
\begin{aligned}
\min_{\{p_g\}_{g\in\srcset}}\quad & \sum_{g\in\srcset} c_g\,p_g \\
\text{s.t.}\quad
& \sum_{g\in\srcset} p_g = D && \text{(power balance)},\\
& 0\le p_g \le \Pmax_g(s,h) && \forall g\in\srcset \quad\text{(capacity)},
\end{aligned}
\label{eq:lp}
\end{equation}
where $p_g$ is the dispatched power of source $g$. When the total firm capacity
is insufficient, the balance constraint is relaxed and the deficit is served as
\emph{load shedding}
\begin{equation}
L(s,h)=\max\!\Bigl(D-\textstyle\sum_{g}\Pmax_g(s,h),\,0\Bigr).
\label{eq:shed}
\end{equation}
A spinning-reserve requirement of fraction $\rho$ is monitored via the reserve
margin
\begin{equation}
\mathrm{RM}(s,h)=\frac{\sum_{g}\bigl(\Pmax_g(s,h)-p_g\bigr)}{D},
\qquad \text{reserve adequate} \iff \mathrm{RM}(s,h)\ge\rho.
\label{eq:reserve}
\end{equation}

\subsection{Merit-order optimality}
Order the sources by ascending marginal cost into the \emph{merit order}
$\pi=(g_{(1)},g_{(2)},\dots)$ with $c_{g_{(1)}}\le c_{g_{(2)}}\le\cdots$.

\begin{theorem}[Greedy merit order is optimal]
\label{thm:merit}
For the single-bus linear program \eqref{eq:lp}, the greedy ``fill-cheapest-first''
rule
\begin{equation}
p_{g_{(k)}}=\min\!\Bigl(D-\textstyle\sum_{j<k}p_{g_{(j)}},\ \Pmax_{g_{(k)}}(s,h)\Bigr)
\label{eq:greedy}
\end{equation}
yields a globally cost-optimal dispatch.
\end{theorem}
\begin{proof}[Proof sketch]
The feasible region of \eqref{eq:lp} is a bounded polytope and the objective is
linear, so an optimum lies at a vertex. Suppose an optimal dispatch
$\{p_g^\star\}$ does not follow \eqref{eq:greedy}: then there exist sources
$a,b$ with $c_a<c_b$, $p_a^\star<\Pmax_a$ and $p_b^\star>0$. Shifting
$\delta=\min(\Pmax_a-p_a^\star,\,p_b^\star)>0$ from $b$ to $a$ preserves the
balance constraint and changes cost by $-(c_b-c_a)\delta<0$, contradicting
optimality. Hence no such pair exists and the optimum coincides with the greedy
fill. \qed
\end{proof}
Theorem~\ref{thm:merit} is what makes the digital twin both fast and provably
optimal: a single linear pass over the merit order replaces a full LP solve at
every hour.

\subsection{Cost and emissions accounting}
With emission factor $e_g$ (\si{tCO_2/MWh}), the hourly operating cost and
carbon output are
\begin{equation}
C(s,h)=\sum_{g}c_g\,p_g,
\qquad
E(s,h)=\sum_{g}e_g\,p_g,
\end{equation}
and seasonal/daily totals are obtained by summing over the 24-hour
representative day.

% =====================================================================
\section{Calibration estimators}
All parameters in \eqref{eq:demand}--\eqref{eq:pmax} are estimated from the
record $\mathcal{D}$. Writing $\mathcal{D}_s$ for the subset of records in
season $s$ and $Q_p(\cdot)$ for the $p$-th percentile,
\begin{align}
\bar{D}_s &= \operatorname{mean}_{\mathcal{D}_s} D, &
\eta_s &= \operatorname{std}_{\mathcal{D}_s} D \,/\, \bar{D}_s, \\
\phi_s(h) &= \frac{\operatorname{mean}\{D : \text{hour}=h,\,\mathcal{D}_s\}}
{\frac1{24}\sum_{h'}\operatorname{mean}\{D:\text{hour}=h',\,\mathcal{D}_s\}}, &
C_g &= \max_{s\in\seasonset} Q_{95}\!\left(\mathcal{D}_s, g\right), \\
\alpha_{g,s} &= \frac{Q_{95}\!\left(\mathcal{D}_s, g\right)}{C_g}\in[0,1], &
\psi_{\text{solar}}(h) &= \frac{\operatorname{mean}\{\text{solar}:\text{hour}=h\}}
{\max_{h'}\operatorname{mean}\{\text{solar}:\text{hour}=h'\}} .
\end{align}
Defining firm capacity as the best seasonal $Q_{95}$ output, and availability as
the seasonal $Q_{95}$ relative to that firm value, lets a single matrix
$\alpha_{g,s}$ absorb fuel-supply, hydrological and outage limits exactly as
they appear in operations. The growth rate $g$ is the slope of a log-linear fit
to the annual peak demand, $\ln\!\big(\max_{\mathcal{D}_y}D\big)=\ln a + (y)\ln(1+g)$.

% =====================================================================
\section{The SAED algorithm}
Algorithm~\ref{alg:saed} states the per-hour dispatch; it is a direct
implementation of \eqref{eq:pmax}--\eqref{eq:reserve} exploiting
Theorem~\ref{thm:merit}.

\begin{algorithm}[htbp]
\caption{Seasonal-Aware Economic Dispatch (one hour)}
\label{alg:saed}
\KwIn{demand $D$; season $s$; hour $h$; calibration
$(\,C_g,\alpha_{g,s},\psi_g,c_g,e_g\,)$; reserve fraction $\rho$}
\KwOut{dispatch $\{p_g\}$; load shed $L$; reserve margin $\mathrm{RM}$;
cost $C$; emissions $E$}
\BlankLine
\ForEach{$g\in\srcset$}{
  $\Pmax_g \leftarrow C_g\,\alpha_{g,s}\,\psi_g(h)$ \tcp*{firm available cap.}
}
$\pi \leftarrow \textsc{Sort}(\srcset \text{ by } c_g \text{ ascending})$
  \tcp*{merit order}
$R \leftarrow D$\;
\ForEach{$g\in\pi$}{
  $p_g \leftarrow \min(R,\ \Pmax_g)$ \tcp*{fill cheapest first}
  $R \leftarrow R-p_g$\;
  \lIf{$R\le 0$}{\textbf{break}}
}
$L \leftarrow \max(R,0)$ \tcp*{unserved energy}
$\mathrm{RM} \leftarrow \big(\sum_g(\Pmax_g-p_g)\big)/D$\;
$C \leftarrow \sum_g c_g p_g$;\quad $E \leftarrow \sum_g e_g p_g$\;
\Return $\{p_g\},L,\mathrm{RM},C,E$\;
\end{algorithm}

% =====================================================================
\section{The digital twin}
The twin (Algorithm~\ref{alg:twin}) wraps SAED in a calibrated, optionally
stochastic, season-by-season replay. For a chosen scenario year it returns a
full 24-hour dispatch trajectory per season, from which all figures and tables
are produced.

\begin{algorithm}[htbp]
\caption{SAED Digital Twin (seasonal replay)}
\label{alg:twin}
\KwIn{record $\mathcal{D}$; scenario year $y$; stochastic flag}
\KwOut{trajectory table $\mathcal{T}$ over $(s,h)$}
\BlankLine
$(\,\bar D_s,\phi_s,\eta_s,C_g,\alpha_{g,s},\psi_g,g,y_0\,)\leftarrow
\textsc{Calibrate}(\mathcal{D})$\;
$\mathcal{T}\leftarrow\varnothing$\;
\ForEach{$s\in\seasonset$}{
  \For{$h\leftarrow 0$ \KwTo $23$}{
    $D \leftarrow \bar D_s\,\phi_s(h)\,(1+g)^{\,y-y_0}$\;
    \lIf{stochastic}{$D\leftarrow D\,(1+\varepsilon),\ \varepsilon\sim
      \mathcal{N}(0,(\tfrac12\eta_s)^2)$}
    $(\{p_g\},L,\mathrm{RM},C,E)\leftarrow\textsc{SAED}(D,s,h,\cdot)$
      \tcp*{Alg.~\ref{alg:saed}}
    append $(s,h,\{p_g\},D,L,\mathrm{RM},C,E)$ to $\mathcal{T}$\;
  }
}
\Return $\mathcal{T}$\;
\end{algorithm}

% =====================================================================
\section{Calibrated parameters}
Tables~\ref{tab:demand}--\ref{tab:source} report the calibrated demand and
techno-economic parameters; the seasonal availability matrix $\alpha_{g,s}$ is
given in Table~\ref{tab:avail} and visualized in
Figure~\ref{fig:heatmap}. The estimated annual peak-demand growth used for the
stress test is $g\approx 7.9\%$ with base year $y_0=2025$.

\input{tables/all_tables.tex}

\begin{figure}[htbp]
\centering
\includegraphics[width=.66\textwidth]{figures/fig4_availability_heatmap.pdf}
\caption{Calibrated seasonal availability matrix $\alpha_{g,s}\in[0,1]$. The
seasonal physics of the grid is visible directly: hydro collapses in winter
($0.41$) and saturates in the monsoon ($1.00$); solar peaks in the pre-monsoon
($1.00$); domestic gas contracts in winter ($0.83$).}
\label{fig:heatmap}
\end{figure}

% =====================================================================
\section{Results: how season changes generation}

\subsection{Seasonal demand and dispatch}
Figure~\ref{fig:demand} shows the calibrated diurnal demand profiles. Winter
carries the lowest base load with a pronounced morning shoulder, whereas the
monsoon and pre-monsoon run $30$--$37\%$ higher with a sharp post-sunset peak.
Feeding these profiles through SAED yields the seasonal dispatch stacks of
Figure~\ref{fig:stacks}: cheap domestic gas forms the base in every season,
coal follows in merit order, and the renewable and import tiers enter according
to the seasonal availability matrix.

\begin{figure}[htbp]
\centering
\includegraphics[width=.62\textwidth]{figures/fig1_demand_profiles.pdf}
\caption{Calibrated seasonal diurnal demand profiles $\bar D_s\phi_s(h)$.}
\label{fig:demand}
\end{figure}

\begin{figure}[htbp]
\centering
\includegraphics[width=\textwidth]{figures/fig2_dispatch_stacks.pdf}
\caption{SAED merit-order dispatch stacks for a representative day in each
season (base year). The dashed line is demand; the colored bands are the
optimal source contributions in merit order.}
\label{fig:stacks}
\end{figure}

\subsection{Validation against the historical mix}
Figure~\ref{fig:valid} compares the twin's seasonal mean dispatch against the
recorded mix, with per-source errors in Table~\ref{tab:valid}. The twin tracks
the gas-dominated base and the seasonal ordering of coal and hydro. The largest
discrepancy is deliberate and economically meaningful: SAED substitutes cheap
gas and coal for the expensive liquid-fuel and import generation the historical
operation relied on, indicating a recoverable efficiency margin rather than a
modelling error.

\begin{figure}[htbp]
\centering
\includegraphics[width=\textwidth]{figures/fig3_validation.pdf}
\caption{Digital-twin validation: SAED dispatch (colored) versus historical
generation (grey) by source and season.}
\label{fig:valid}
\end{figure}

\subsection{Cost, carbon and the growth stress test}
Seasonal operating cost and emissions (Figure~\ref{fig:cost}) scale with demand
and peak in the monsoon. The demand-growth stress test
(Figure~\ref{fig:stress}) propagates the calibrated $7.9\%$ annual growth
forward: liquid-fuel peaking energy rises fastest in the monsoon, and the
monsoon peak-hour reserve margin is the first to fall below the $\rho=8\%$
planning threshold (around 2029--2030), identifying the monsoon as the binding
season for capacity expansion.

\begin{figure}[htbp]
\centering
\includegraphics[width=\textwidth]{figures/fig5_cost_emissions.pdf}
\caption{Seasonal daily operating cost (a) and CO$_2$ emissions (b) from the
optimal SAED dispatch.}
\label{fig:cost}
\end{figure}

\begin{figure}[htbp]
\centering
\includegraphics[width=\textwidth]{figures/fig6_stress_test.pdf}
\caption{Demand-growth stress test by season: (a) liquid-fuel peaking energy and
(b) peak-hour reserve margin. The monsoon crosses the $8\%$ reserve floor
first.}
\label{fig:stress}
\end{figure}

% =====================================================================
\section{Discussion and conclusion}
The SAED algorithm reduces seasonal dispatch to a calibrated demand model, a
seasonal availability matrix, and a provably optimal merit-order pass
(Theorem~\ref{thm:merit}). Driven by ten years of PGCB data, the digital twin
reproduces the grid's seasonal physics and exposes three actionable findings:
(i) the historical reliance on liquid-fuel peaking is largely economically
avoidable under optimal dispatch; (ii) seasonal generation differences are
governed chiefly by hydro and solar availability and by the monsoon demand peak,
quantified by the $\alpha_{g,s}$ matrix; and (iii) under the historical growth
trend the monsoon reserve margin is the first to breach the planning floor,
making monsoon firm capacity the priority for expansion. The twin is fully
reproducible from the accompanying calibration file and source code.

\end{document}
